



















\subsection{置换展开}
前一小节定义了：满足四个条件的函数就是行列式，并证明了对每个 $n$，
$\nbyn{n}$ 行列式函数至多有一个。
剩下的工作就是证明对每个 $n$ 都存在这样的函数。

但是，我们很容易计算行列式：
用高斯消元法，记录行对换带来的符号变化，
最后沿对角线相乘。
行列式怎么可能不存在呢？

困难在于证明该计算给出的是良定义的\Dash 即唯一的\Dash 结果。
%<*DifferentGaussMethodReductions>
考虑对同一个矩阵进行如下两次高斯消元法化简，
第一次不包含行对换
\begin{equation*}
  \begin{mat}[r]
    1  &2  \\
    3  &4
  \end{mat}
  \grstep{-3\rho_1+\rho_2}
  \begin{mat}[r]
    1  &2  \\
    0  &-2
  \end{mat}
\end{equation*}
第二次包含一次行对换。
\begin{equation*}
  \begin{mat}[r]
    1  &2  \\
    3  &4
  \end{mat}
  \grstep{\rho_1\leftrightarrow\rho_2}
  \begin{mat}[r]
    3  &4  \\
    1  &2
  \end{mat}
  \grstep{-(1/3)\rho_1+\rho_2}
  \begin{mat}[r]
    3  &4  \\
    0  &2/3
  \end{mat}
\end{equation*}
两者都得到行列式 $-2$，
因为在第二种化简中我们记下了：行对换会改变
沿对角线相乘所得结果的符号。
%</DifferentGaussMethodReductions>
我们可以按两种方式进行计算，这一事实引出了两种方式可能给出不同答案的问题。
% To illustrate how a computation that is like the ones that we are doing could 
% fail to be well-defined, 
% suppose that \nearbydefinition{def:Det} did not include condition~(2). 
% That is, suppose that we instead tried to define determinants so that
% the value would not change on a row swap.
% Then first reduction above would
% yield $-2$ while the second would yield $+2$.
% We could still do computations but they wouldn't give consistent outcomes\Dash
% there is no 
% function that satisfies conditions (1), (3), (4), and also this
% altered second condition. 
% Of course, observing that \nearbydefinition{def:Det} does the right thing
% with these two reductions of the above matrix is not enough.
也就是说，我们给出的这种计算行列式值的方法，并没有明显排除如下可能性：比如说，
某个 $\nbyn{7}$~矩阵的两次化简可能导出
不同的行列式值。
若是那样，我们就不再有函数了，因为函数的定义要求每个输入恰好对应唯一的输出。
本节余下的部分将证明，定义~\nearbydefinition{def:Det}
永远不会导致冲突。

% \begin{remark}
% The natural way approach to this would be 
% to define the determinant function with:~``The value of the
% function is the result of doing Gauss's Method,
% keeping track of row swaps, and finishing by multiplying down the diagonal''. 
% (Since Gauss's Method allows for some variation, as we saw above, 
% we would have to fix an explicit algorithm.)
% Then we would be done if we verified that 
% this way algorithm satisfies the four properties.
% However, how to verify this is not evident.
% So the development below will not proceed in this way.
% \end{remark}

为此，我们将定义另一种求行列式值的方法。
（这种替代方法在实践中用处不大，因为它
速度慢。
但它对理论非常有用。）
关键的想法是，
\nearbydefinition{def:Det} 的条件~(3)
表明行列式函数不是线性的。

\begin{example}
利用条件~(3)，标量可以从每一行中分别提出来，
\begin{equation*}
   \begin{vmat}[r]
       4  &2  \\
      -2  &6
     \end{vmat}
  =2\cdot\begin{vmat}[r]
       2  &1  \\
      -2  &6
     \end{vmat}
  =4\cdot\begin{vmat}[r]
       2  &1  \\
      -1  &3
     \end{vmat}
\end{equation*}
而不是从整个矩阵一次提出。
于是，当
\begin{equation*}
   A=\begin{mat}[r]
       2  &1  \\
      -1  &3
     \end{mat}
\end{equation*} 
时，\( \det(2A) \neq 2\cdot\det(A) \) 
（相反，\( \det(2A) =  4\cdot\det(A) \)）。
\end{example}

由于标量是逐行提出的，我们可以猜测
行列式是逐行线性的。

\begin{definition} \label{def:multilinear}
%<*df:Multilinear>
设 \( V \) 是一个向量空间。
映射 \( \map{f}{V^n}{\Re} \) 称为
\definend{多重线性的}\index{function!multilinear}\index{multilinear}，
若  
\begin{enumerate}
  \item 
    $
      f(\vec{\rho}_1,\dots,\vec{v}+\vec{w},
      \ldots,\vec{\rho}_n)
      =f(\vec{\rho}_1,\dots,\vec{v},\dots,\vec{\rho}_n)
      +f(\vec{\rho}_1,\dots,\vec{w},\dots,\vec{\rho}_n)
    $
  \item 
    $
      f(\vec{\rho}_1,\dots,k\vec{v},\dots,\vec{\rho}_n)
      =k\cdot f(\vec{\rho}_1,\dots,\vec{v},\dots,\vec{\rho}_n)
    $
\end{enumerate}
其中 \( \vec{v}, \vec{w}\in V \)，\( k\in\Re \)。
%</df:Multilinear>
\end{definition}

\begin{lemma}  \label{lem:DetsMultilinear}
%<*lm:DetsMultilinear>
行列式是多重线性的。
%</lm:DetsMultilinear>
\end{lemma}

\begin{proof}
%<*pf:DetsMultilinear0>
这里的性质~(2) 恰是 \nearbydefinition{def:Det} 的条件~(3)，
所以我们只需验证性质~(1)。
%</pf:DetsMultilinear0>
% \begin{equation*}
%   \det(\vec{\rho}_1,\dots,\vec{v}+\vec{w},
%       \dots,\vec{\rho}_n)
%   =\det(\vec{\rho}_1,\dots,\vec{v},\dots,\vec{\rho}_n)
%   +\det(\vec{\rho}_1,\dots,\vec{w},\dots,\vec{\rho}_n)
% \end{equation*}

%<*pf:DetsMultilinear1>
分两种情况。
若其余各行构成的集合
\(
  \set{\vec{\rho}_1,\dots,\vec{\rho}_{i-1},\vec{\rho}_{i+1},
       \dots,\vec{\rho}_n} 
\)
是线性相关的，则这三个矩阵都是奇异的，从而三个行列式都为零，等式平凡地成立。
%</pf:DetsMultilinear1>

%<*pf:DetsMultilinear2>
因此假设其余各行构成的集合是线性无关的。
我们可以再添加一个向量
$\sequence{\vec{\rho}_1,\dots,\vec{\rho}_{i-1},\vec{\beta},
               \vec{\rho}_{i+1},\dots,\vec{\rho}_n}$
来作出一个基。
关于这个基表示 $\vec{v}$ 和 $\vec{w}$
\begin{align*}
  \vec{v} &=v_1\vec{\rho}_1+\dots+v_{i-1}\vec{\rho}_{i-1}+v_i\vec{\beta}
            +v_{i+1}\vec{\rho}_{i+1}+\dots+v_n\vec{\rho}_n                \\
  \vec{w} &= w_1\vec{\rho}_1+\dots+w_{i-1}\vec{\rho}_{i-1}+w_i\vec{\beta}
            +w_{i+1}\vec{\rho}_{i+1}+\dots+w_n\vec{\rho}_n
\end{align*}
再相加。
\begin{equation*}
  \vec{v}+\vec{w}
  =
  (v_1+w_1)\vec{\rho}_1+\dots+(v_i+w_i)\vec{\beta}
            +\dots+(v_n+w_n)\vec{\rho}_n 
\end{equation*}
考虑 (1) 的左边并展开 $\vec{v}+\vec{w}$。
% $\det (\vec{\rho}_1,\dots,\vec{v}+\vec{w},\dots,\vec{\rho}_n)$.
\begin{equation*}
    \det (\vec{\rho}_1,\dots,\,(v_1+w_1)\vec{\rho}_1+\dots+(v_i+w_i)\vec{\beta}
            +\dots+(v_n+w_n)\vec{\rho}_n,\,\dots,\vec{\rho}_n) 
   \tag{$*$}
\end{equation*}
% Gauss's Method enables us to eliminate the terms for $\vec{\rho}_1$,
% $\vec{\rho}_2$, etc., from the expanded expression.
% the $(v_1+w_1)\vec{\rho}_1+\dots+(v_i+w_i)\vec{\beta}
%             +\dots+(v_n+w_n)\vec{\rho}_n$ expression.
由行列式定义的条件~(1)，($*$) % the determinant
% $\det(\vec{\rho}_1,\dots,\vec{v}+\vec{w},\dots,\vec{\rho}_n)$
的值在下述操作下不变：
将 $-(v_1+w_1)\vec{\rho}_1$ 倍加到第 $i$ 行 $\vec{v}+\vec{w}$ 上。
第 $i$ 行变成这样。
\begin{equation*}
  \vec{v}+\vec{w}-(v_1+w_1)\vec{\rho}_1
  =
  (v_2+w_2)\vec{\rho}_2+\cdots+(v_i+w_i)\vec{\beta}
            +\dots+(v_n+w_n)\vec{\rho}_n 
\end{equation*}
%</pf:DetsMultilinear2>
%<*pf:DetsMultilinear3>
接着倍加 $-(v_2+w_2)\vec{\rho}_2$ 等等，以消去
来自其他各行的所有项。
应用行列式定义中的条件~(3)。
\begin{multline*}
  \det (\vec{\rho}_1,\dots,\vec{v}+\vec{w},\dots,\vec{\rho}_n)          \\
  \begin{aligned}
    &=\det (\vec{\rho}_1,\dots,(v_i+w_i)\cdot\vec{\beta},\dots,\vec{\rho}_n) \\
    &=(v_i+w_i)\cdot\det (\vec{\rho}_1,\dots,\vec{\beta},\dots,\vec{\rho}_n) \\
    &=v_i\cdot \det (\vec{\rho}_1,\dots,\vec{\beta},\dots,\vec{\rho}_n)  
     +w_i\cdot \det (\vec{\rho}_1,\dots,\vec{\beta},\dots,\vec{\rho}_n)
  \end{aligned}
\end{multline*}
现在这是两个行列式的和。
为完成证明，把 $v_i$ 与 $w_i$ 重新移回到
$\vec{\beta}$ 的前面，
并再次使用行倍加，
这一次用来按基重构 $\vec{v}$
和 $\vec{w}$ 的表达式。
也就是说，从下述操作开始：
将 $v_1\vec{\rho}_1$ 倍加到 $v_i\vec{\beta}$，
将 $w_1\vec{\rho}_1$ 倍加到 $w_i\vec{\rho}_1$，等等，
从而得到 $\vec{v}$ 和 $\vec{w}$ 的展开式。
%</pf:DetsMultilinear3>
\end{proof}

多重线性使我们能把一个行列式展开成一些行列式的和，
其中每一个都涉及一个简单的矩阵。

\begin{example}
利用多重线性的性质~(1)
拆分第一行
\begin{equation*}
  \begin{vmat}[r]
     2  &1  \\
     4  &3
  \end{vmat}
  =
  \begin{vmat}[r]
     2  &0  \\
     4  &3
  \end{vmat}
  +
  \begin{vmat}[r]
     0  &1  \\
     4  &3
  \end{vmat}
\end{equation*}
然后再用 (1) 沿第二行拆分每一个。
\begin{equation*}
  =\begin{vmat}[r]
     2  &0  \\
     4  &0
  \end{vmat}
  +
  \begin{vmat}[r]
     2  &0  \\
     0  &3
  \end{vmat}
  +
  \begin{vmat}[r]
     0  &1  \\
     4  &0
  \end{vmat}
  +
  \begin{vmat}[r]
     0  &1  \\
     0  &3
  \end{vmat}
\end{equation*}
结果是四个行列式。
在这四个行列式的每一行中，
都只有来自原矩阵的一个元素。
\end{example}

\begin{example}
同样地，
一个 \( \nbyn{3} \) 行列式分离成
许多更简单行列式的和。
沿第一行拆分产生三个行列式
（我们高亮了 $1,3$ 位置的零，使它在视觉上
与作为拆分一部分出现的零区分开）。
\begin{equation*}
  \begin{vmat}[r]
     2              &1  &-1  \\
     4              &3  &\highlight{0}  \\
     2              &1  &5
  \end{vmat}
  =
  \begin{vmat}[r]
     2              &0  &0   \\
     4              &3  &\highlight{0}  \\
     2              &1  &5
  \end{vmat}
  +
  \begin{vmat}[r]
     0              &1  &0   \\
     4              &3  &\highlight{0}   \\
     2              &1  &5
  \end{vmat}
  +
  \begin{vmat}[r]
     0              &0  &-1  \\
     4              &3  &\highlight{0}  \\
     2  &1  &5
  \end{vmat}                       
\end{equation*}
反过来，上面每一个都沿第二行拆分成三个。
然后这九个每一个再沿第三行拆分成三个。
结果是二十七个行列式，
使得每一行都只包含来自初始矩阵的一个元素。
\begin{equation*}
  =
  \begin{vmat}[r]
     2              &0  &0   \\
     4              &0  &0   \\
     2              &0  &0
  \end{vmat}
  +
  \begin{vmat}[r]
     2  &0  &0   \\
     4  &0  &0   \\
     0  &1  &0
  \end{vmat}
  +
  \begin{vmat}[r]
     2  &0  &0   \\
     4  &0  &0   \\
     0  &0  &5
  \end{vmat}
  +
  \begin{vmat}[r]
     2  &0  &0   \\
     0  &3  &0   \\
     2  &0  &0
  \end{vmat}
  +\dots+
  \begin{vmat}[r]
     0  &0  &-1  \\
     0  &0  &\highlight{0}  \\
     0  &0  &5
  \end{vmat}
\end{equation*}
\end{example}

于是多重线性会将一个 \( \nbyn{n} \) 行列式展开成
\( n^n \) 个行列式的和，其中
每个行列式的每一行都只包含来自
初始矩阵的一个元素。

在这个展开中，尽管项很多，
但大多数的行列式值为零。

\begin{example} \label{ex:SamplePermExp}
在来自前面展开式的这些例子中，
原矩阵的两个元素位于同一列。
\begin{equation*}
  \begin{vmat}[r]
     2               &0  &0   \\
     4               &0  &0  \\
     0               &1  &0
  \end{vmat}
  \qquad
  \begin{vmat}[r]
     0               &0  &-1  \\
     0               &3  &0  \\
     0               &0  &5
  \end{vmat}
  \qquad\
  \begin{vmat}[r]
     0               &1  &0   \\
     0               &0  &\highlight{0}  \\
     0               &0  &5
  \end{vmat}
\end{equation*}
例如，
在第一个矩阵中，$2$ 和 $4$ 都来自原矩阵的第一列。
在第二个矩阵中，$-1$ 和~$5$ 都来自第三列。
而在第三个矩阵中，$0$ 和~$5$ 都来自第三列。
任何这样的矩阵都是奇异的，因为
其中一行是另一行的倍数。
因此，根据 \nearbylemma{le:IdenRowsDetZero}，任何这样的行列式都为零。

有了这一观察，
上面把 \( \nbyn{3} \) 行列式展开成
二十七个行列式之和的做法就化简为这六个行列式之和：
其中来自原矩阵的元素每行一个，
且每列也只有一个。
\begin{align*}
  \begin{vmat}[r]
     2  &1  &-1  \\
     4  &3  &\highlight{0}  \\
     2  &1  &5
  \end{vmat}
  &=\begin{vmat}[r]
     2  &0  &0   \\
     0  &3  &0   \\
     0  &0  &5
  \end{vmat}
  +
   \begin{vmat}[r]
     2  &0  &0   \\
     0  &0  &\highlight{0}   \\
     0  &1  &0
  \end{vmat}                      \\
  &\quad\hbox{}+\begin{vmat}[r]
     0  &1  &0   \\
     4  &0  &0   \\
     0  &0  &5
  \end{vmat}
  +
   \begin{vmat}[r]
     0  &1  &0   \\
     0  &0  &\highlight{0}   \\
     2  &0  &0
  \end{vmat}                      \\
  &\quad\hbox{}+\begin{vmat}[r]
     0  &0  &-1  \\
     4  &0  &0   \\
     0  &1  &0
  \end{vmat}
  +
   \begin{vmat}[r]
     0  &0  &-1  \\
     0  &3  &0    \\
     2  &0  &0
  \end{vmat}                      
\end{align*}
在那个展开式中我们可以把标量提出来。
\begin{align*}
  &=(2)(3)(5)\begin{vmat}[r]
               1  &0  &0  \\
               0  &1  &0  \\
               0  &0  &1
            \end{vmat}
  +(2)(\highlight{0})(1)\begin{vmat}[r]
               1  &0  &0  \\
               0  &0  &1  \\
               0  &1  &0
            \end{vmat}                     \\
  &\quad\hbox{}+(1)(4)(5)\begin{vmat}[r]
               0  &1  &0  \\
               1  &0  &0  \\
               0  &0  &1
            \end{vmat}
  +(1)(\highlight{0})(2)\begin{vmat}[r]
               0  &1  &0  \\
               0  &0  &1  \\
               1  &0  &0
            \end{vmat}                       \\
  &\quad\hbox{}+(-1)(4)(1)\begin{vmat}[r]
               0  &0  &1  \\
               1  &0  &0  \\
               0  &1  &0
            \end{vmat}
  +(-1)(3)(2)\begin{vmat}[r]
               0  &0  &1  \\
               0  &1  &0  \\
               1  &0  &0
            \end{vmat}                       
\end{align*}
为完成计算，通过行对换把那六个行列式
化为单位矩阵，
并记录符号的变化。
\begin{align*}
  &=30\cdot (+1)+0\cdot (-1)  \\
  &\quad\hbox{}+20\cdot (-1)+0\cdot (+1) \\
  &\quad \hbox{}-4\cdot (+1)-6\cdot (-1)=12
\end{align*}
\end{example}

这个例子体现了本小节新的计算方案。
多重线性把一个行列式展开成
许多单独的行列式，每个行列式的每一行有来自
原矩阵的一个元素。
其中大多数都
有一行是另一行的倍数，
所以我们把它们略去。
剩下的是那些每行每列各有一个原矩阵元素的行列式。
提出标量进一步把我们必须计算的行列式
化简为每行每列恰有一个元素的、所有元素都是 $1$ 的矩阵。

%<*NotationForPermutationMatrices>
回忆定义~Three.IV.\ref{df:PermutationMatrix}：
一个 \definend{置换矩阵}\index{permutation matrix}%
\index{matrix!permutation}
是方阵，其元素除每行每列恰有一个 $1$ 外都是 $0$。
现在我们引入置换矩阵的一种记号。
%</NotationForPermutationMatrices>

\begin{definition}  \label{df:permutation}
%<*df:permutation>
一个 \definend{\( n \)-置换}\index{permutation}
是前~$n$ 个正整数上的一个函数
$\map{\phi}{\set{1,\ldots,n}}{\set{1,\ldots,n}}$，
它是单射且满射的。
%</df:permutation>
\end{definition}

在一个置换中，每个数
$1$, \ldots, $n$ 都作为且仅作为一个输入的输出出现。
我们可以把置换记为一个序列
$\phi=\sequence{\phi(1),\phi(2),\ldots,\phi(n)}$。

\begin{example} \label{ex:AllTwoThreePerms}
%<*ex:AllTwoPerms>
$2$-置换是函数
$\map{\phi_1}{\set{1,2}}{\set{1,2}}$
由 $\phi_1(1)=1$，$\phi_1(2)=2$ 给出，
以及
$\map{\phi_2}{\set{1,2}}{\set{1,2}}$
由 $\phi_2(1)=2$，$\phi_2(2)=1$ 给出。
序列记号更简洁：
\( \phi_1=\sequence{1,2} \) 与 \( \phi_2=\sequence{2,1} \)。
%</ex:AllTwoPerms>
\end{example}

\begin{example} \label{ex:AllThreePerms}
%<*ex:AllThreePerms>
在序列记号下，$3$-置换是
\( \phi_1=\sequence{1,2,3} \)、
\( \phi_2=\sequence{1,3,2} \)、
\( \phi_3=\sequence{2,1,3} \)、
\( \phi_4=\sequence{2,3,1} \)、
\( \phi_5=\sequence{3,1,2} \)，以及
\( \phi_6=\sequence{3,2,1} \)。
%</ex:AllThreePerms>
\end{example}

%<*AssociatePermutationWithPermutationMatrix>
我们把除第 $j$ 个元素处为 $1$ 外全为 $0$ 的行向量记为 $\iota_j$，
于是宽度为四的 $\iota_2$
是~$\rowvec{0 &1 &0 &0}$。
现在我们的置换矩阵记号是：
对任意 $\phi=\sequence{\phi(1),\ldots,\phi(n)}$，
与之相伴的矩阵的行依次为 $\iota_{\phi(1)}$, \ldots, $\iota_{\phi(n)}$。
%</AssociatePermutationWithPermutationMatrix>
例如，与 $4$-置换 \( \phi=\sequence{3,2,1,4} \) 相伴的
是其行依次为相应 $\iota$ 的矩阵。
\begin{equation*}
   P_{\phi}=
   \begin{mat}
      \iota_{3} \\
      \iota_{2} \\
      \iota_{1} \\
      \iota_{4} 
   \end{mat}=
  \begin{mat}[r]
     0  &0  &1  &0  \\
     0  &1  &0  &0  \\
     1  &0  &0  &0  \\
     0  &0  &0  &1
  \end{mat}
\end{equation*}

\begin{example} \label{ex:TwoThreePermsAndMatrices}
这些是 \nearbyexample{ex:AllTwoThreePerms} 中列出的
$2$-置换的置换矩阵。
\begin{equation*}
  P_{\phi_1}
  =\begin{mat}
      \iota_1 \\
      \iota_2 
   \end{mat}
  =\begin{mat}[r]
    1  &0         \\
    0  &1   
  \end{mat}
  \qquad
  P_{\phi_2}
   =\begin{mat}
      \iota_2 \\
      \iota_1 
   \end{mat}
   =\begin{mat}[r]
    0  &1         \\
    1  &0   
  \end{mat}
\end{equation*}
例如，$P_{\phi_2}$ 的第一行是
$\iota_{\phi_2(1)}=\iota_2$，第二行是 $\iota_{\phi_2(2)}=\iota_1$。
\end{example}

\begin{example}
考虑 $3$-置换
\( \phi_5=\sequence{3,1,2} \)。
置换矩阵  
$P_{\phi_5}$ 的行依次为 $\iota_{\phi_5(1)}=\iota_3$、
$\iota_{\phi_5(2)}=\iota_1$ 与 $\iota_{\phi_5(3)}=\iota_2$。\begin{equation*}
  % P_{\phi_2}
  %  =\begin{mat}
  %     \iota_1 \\
  %     \iota_3 \\
  %     \iota_2 
  %  \end{mat}
  % =\begin{mat}[r]
  %    1      &0        &0        \\
  %    0      &0        &1        \\
  %    0      &1        &0
  % \end{mat}
  % \qquad
  P_{\phi_5}
   % =\begin{mat}
   %    \iota_3 \\
   %    \iota_1 \\
   %    \iota_2 
   % \end{mat}
  =\begin{mat}[r]
     0      &0        &1        \\
     1      &0        &0        \\
     0      &1        &0
  \end{mat}
\end{equation*}
\end{example}

\begin{definition} \label{df:PermutationExpansion}
%<*df:PermutationExpansion>
行列式的 \definend{置换展开}\index{determinant!permutation expansion}%
\index{permutation expansion}
是
\begin{equation*}
   \begin{vmat}
      t_{1,1}  &t_{1,2}  &\ldots  &t_{1,n}  \\
      t_{2,1}  &t_{2,2}  &\ldots  &t_{2,n}  \\
               &\vdots                      \\
      t_{n,1}  &t_{n,2}  &\ldots  &t_{n,n}
   \end{vmat}
   =
   \begin{array}[t]{l}
      t_{1,\phi_1(1)}t_{2,\phi_1(2)}\cdots
           t_{n,\phi_1(n)}\deter{P_{\phi_1}}       \\[.5ex]
      \>\hbox{}+t_{1,\phi_2(1)}t_{2,\phi_2(2)}\cdots
           t_{n,\phi_2(n)}\deter{P_{\phi_2}}       \\
      \>\vdotswithin{+}                              \\
      \>\hbox{}+t_{1,\phi_k(1)}t_{2,\phi_k(2)}\cdots
           t_{n,\phi_k(n)}\deter{P_{\phi_k}} 
   \end{array}
\end{equation*}
其中 \( \phi_1,\ldots,\phi_k \) 是所有的 \( n \)-置换。
%</df:PermutationExpansion>
\end{definition}
%<*SummationForPermutationExpansion>
我们可以把这个公式改述为
\definend{求和
记号}\index{summation notation, for permutation expansion}
\begin{equation*}
  \deter{T}=\!\!
  \sum_{\text{置换\ }\phi}\!\!\!\!
     t_{1,\phi(1)}t_{2,\phi(2)}\cdots t_{n,\phi(n)}
                                 \deter{P_{\phi}}
\end{equation*}
读作，
"对所有的置换 \( \phi \) 求和，项形如
\( t_{1,\phi(1)}t_{2,\phi(2)}\cdots t_{n,\phi(n)} \deter{P_{\phi}} \)。"
%</SummationForPermutationExpansion>


\begin{example}
熟悉的 $\nbyn{2}$ 行列式公式可由上面推出
\begin{align*}
  \begin{vmat}
    t_{1,1}  &t_{1,2} \\
    t_{2,1}  &t_{2,2}
  \end{vmat}
  &=
  t_{1,1}t_{2,2}\cdot\deter{P_{\phi_1}}
  +
  t_{1,2}t_{2,1}\cdot\deter{P_{\phi_2}}      \\[-1.5ex]     
  &=
  t_{1,1}t_{2,2}\cdot\begin{vmat}[r]
           1  &0 \\
           0  &1
         \end{vmat}
  +
  t_{1,2}t_{2,1}\cdot\begin{vmat}[r]
           0  &1 \\
           1  &0
         \end{vmat}               \\
  &=t_{1,1}t_{2,2}-t_{1,2}t_{2,1}
\end{align*}
$\nbyn{3}$ 公式也是如此。
\begin{align*}
  \begin{vmat}
    t_{1,1}  &t_{1,2}  &t_{1,3} \\
    t_{2,1}  &t_{2,2}  &t_{2,3} \\
    t_{3,1}  &t_{3,2}  &t_{3,3} 
  \end{vmat}
  &=
  \begin{aligned}[t]
    &t_{1,1}t_{2,2}t_{3,3}\deter{P_{\phi_1}}
     +t_{1,1}t_{2,3}t_{3,2}\deter{P_{\phi_2}}
     +t_{1,2}t_{2,1}t_{3,3}\deter{P_{\phi_3}} \\
    &\hspace*{0em}
     +t_{1,2}t_{2,3}t_{3,1}\deter{P_{\phi_4}}
     +t_{1,3}t_{2,1}t_{3,2}\deter{P_{\phi_5}}
     +t_{1,3}t_{2,2}t_{3,1}\deter{P_{\phi_6}}
  \end{aligned}                                      \\
  &=
  \begin{aligned}[t]
    &t_{1,1}t_{2,2}t_{3,3}
     -t_{1,1}t_{2,3}t_{3,2}
     -t_{1,2}t_{2,1}t_{3,3}  \\
    &\hspace*{0em}
     +t_{1,2}t_{2,3}t_{3,1}
     +t_{1,3}t_{2,1}t_{3,2}
     -t_{1,3}t_{2,2}t_{3,1}
  \end{aligned}
\end{align*}
\end{example}

用置换展开计算行列式通常比用高斯消元法
耗时更长。
不过，我们将用它来证明
行列式函数存在。
这个证明很长，所以我们在这里只陈述结果，
而把证明留到下一小节。

\begin{theorem} \label{th:DetsExist}
\index{determinant!exists} 
%<*th:DetsExist>
对每个 $n$，都存在一个 $\nbyn{n}$ 行列式函数。
%</th:DetsExist>
\end{theorem}

下一小节还将给出
下一个结果的证明（把它们放在一起是因为两个证明有重叠）。

\begin{theorem}  \label{th:DeterminantOfAMatrixEqualsDeterminantOfTranspose}
\index{transpose!determinant}
%<*th:DeterminantOfAMatrixEqualsDeterminantOfTranspose>
一个矩阵的行列式等于其转置的行列式。
%</th:DeterminantOfAMatrixEqualsDeterminantOfTranspose>
\end{theorem}

由于这个定理，
虽然迄今为止我们都是按行陈述行列式的结果，
但所有这些结果按列也同样成立。

\begin{corollary} \label{cor:ColSwapChgSign} \label{cor:DetsMultiInCols}
%<*co:ColSwapChgSign>
  有两列相等的矩阵是奇异的。
  列对换改变行列式的符号。
  行列式关于其列是多重线性的。
%</co:ColSwapChgSign>
\end{corollary}

\begin{proof}
%<*pf:ColSwapChgSign>
对于第一个命题，
转置该矩阵得到一个具有相同行列式的矩阵，
且它有两行相等，从而行列式为零。
其余两个命题按同样的方式证明。
%</pf:ColSwapChgSign>
\end{proof}

我们以一段总结来结束本小节：
行列式函数存在、唯一，并且我们知道如何计算它们。
至于行列式究竟是关于什么的，下面几行文字
\cite{Kemp} 或许有助于我们记住它。
\begin{verse} \small
行列式为零，         \\*
\ 解：解无穷多或无解。   \\*
行列式非零，         \\*
\ 解：恰有一解。
\end{verse}




